\begin{figure}[htbp]
\centering
\begin{adjustbox}{max width=0.98\linewidth,max totalheight=0.72\textheight}
\begin{tikzpicture}[
    node distance=2cm,
    auto,
    block/.style={
        rectangle,
        draw,
        thick,
        fill=white,
        text width=3cm,
        align=center,
        minimum height=3cm
    },
    line/.style={
        draw, -latex', thick
    },
    interface/.style={
        rectangle,
        draw,
        fill=gray!20,
        minimum height=4cm,
        minimum width=1cm,
        align=center,
        rotate=0
    }
]

% Nodes
\node [block, fill=blue!10] (geo) {\textbf{几何核}\\(TPU-G)\\ \footnotesize{连续流形 $\mathcal{M}$}\\ \footnotesize{模拟/光子}\\ \footnotesize{直觉/本理论框架}};
\node [interface, right of=geo, node distance=4cm] (wpi) {\textbf{W}\\ \textbf{P}\\ \textbf{I}};
\node [block, fill=red!10, right of=wpi, node distance=4cm] (alg) {\textbf{代数核}\\(Logic Core)\\ \footnotesize{离散符号 $\mathcal{S}$}\\ \footnotesize{数字逻辑}\\ \footnotesize{推理/Turing}};

% Arrows
\draw [->, thick] (geo.20) -- node [above, font=\scriptsize] {波 $\Psi$ (模拟)} (wpi.160);
\draw [->, thick] (wpi.20) -- node [above, font=\scriptsize] {粒子语义子 (数字)} (alg.160);

\draw [->, thick] (alg.200) -- node [below, font=\scriptsize] {指令/控制 (数字)} (wpi.340);
\draw [->, thick] (wpi.200) -- node [below, font=\scriptsize] {势能/场 $\mathbf{\Gamma}$ (模拟)} (geo.340);

\node [below of=wpi, node distance=3cm] {波粒二象性接口};

\end{tikzpicture}
\end{adjustbox}
\caption{双院制芯片架构：几何核与代数核的物理耦合}
\label{fig:bicameral_chip}
\end{figure}
